% -*- root: ../thesis.tex -*-
%!TEX root = ../thesis.tex
% ******************************* Thesis Chapter PPT ****************************


% ----------------------- paths to graphics ------------------------

% change according to folder and file names
\graphicspath{{chapter_PPT/figures/}}
\
% ----------------------- contents from here ------------------------

% Could also remove section heading here
\section{Context within thesis}

This chapter introduces the parallel pathway theory (PPT) as a circuit mechanism for systems memory consolidation. It's central hypothesis is that feedforward memory transfer between brain regions can arise form Hebbian plasticity in parallel synaptic pathways---two features that are ubiquitous in the brain. The content of this chapter has been published in a peer-reviewed journal article \cite{Remme2021}.

%%% ABSTRACT

\todo[inline]{Do I need to add references to the abstract here? Needs consistency across the thesis!}

Systems memory consolidation involves the transfer of memories across brain regions and the transformation of memory content. For example, declarative memories that transiently depend on the hippocampal formation are transformed into long-term memory traces in neo-cortical networks, and procedural memories are transformed within cortico-striatal networks. These consolidation processes are thought to rely on replay and repetition of recently acquired memories, but the cellular and network mechanisms that mediate the changes of memories are poorly understood. Here, we suggest that systems memory consolidation could arise from Hebbian plasticity in networks with parallel synaptic pathways --- two ubiquitous features of neural circuits in the brain. We explore this hypothesis in mathematical analyses and computational models to illustrate how memories are transferred across circuits and how this could mediate the semantization of episodic memories. These modelling results are in quantitative agreement with lesion studies in rodents. A hierarchical iteration of the mechanism yields power-law forgetting -- as observed in psychophysical studies in humans, and it yields gradients of episodic and semantic memory strength across the hierarchy. Moreover, the results suggest that memory replay is only necessary for the transfer of episodic but not semantic memories. The predicted circuit mechanism thus bridges spatial scales from single cells to cortical areas, time scales from milliseconds to years, and connects the transfer of memories to their transformation into semantic knowledge.

\todo[inline]{Removed the linear approximation sentence}


%\section{Contribution}
%\ldots
%
%\section{Remme 2021 paper}

\section{Chapter summary}
This chapter introduced the parallel pathway theory (PPT) as a general circuit mechanism for feedforward memory transfer. This theory builds the basis for the following chapters. Part II of this thesis will explore how episodic memories can transform to semantic ones during consolidation via the PPT (chapter \ref{ch:ppt-semantization}) and how different levels of semantization can implement Bayesian change point detection in dynamic environments (chapter \ref{ch:change-point-detection}). Part III of this thesis will illustrate an implementation of the PPT in the mushroom body of Drosophila and illustrate how consolidation in the mushroom body can be seen as hidden feedforward structure in a recurrent network network (chapter \ref{ch:drosophila}) and how the emergent dynamics can be generalized to brain-wide engram dynamics under systems consolidation and representational drift (chapter \ref{ch:engram-reorganization}).



% ---------------------------------------------------------------------------
% ----------------------- end of thesis sub-document ------------------------
% ---------------------------------------------------------------------------
